\section{Distributivité de la multiplication par rapport à l'addition et la soustraction}

\newcommand{\myk}{\textcolor{Bleu}{k}}
\newcommand{\mya}{\textcolor{Rouge}{a}}
\newcommand{\myb}{\textcolor{Vert}{b}}
\begin{prop}
    Pour tous les nombres $\myk$, $\mya$ et $\myb$, on a :
    $$\myk\times(\mya+\myb)=\answer{\myk\times \mya + \myk \times \myb}\text{\hspace{1.5cm} et \hspace{1.5cm}} \myk\times (\mya-\myb) = \answer{\myk\times \mya - \myk\times \myb}$$
    On dit qu'on a \bb{développé} l'expression $\myk\times(\mya+\myb)$.
\end{prop}
\begin{exe}
    \bb{Calculer} $4\times(2+3)$ \bb{en utilisant la règle de distributivité}. 
    \begin{align*}
        4\times(2+3) &= \answer{4\times2} + \answer{4\times3} \\
                     &= \answer{8+12} \\
                     &= \answer{20}
    \end{align*}

    \bb{Développer} puis \bb{calculer} $10\times(6-4)$.
    \begin{align*}
        10\times(6-4) &= \answer{10\times6} - \answer{10\times4}\\
                      &= \answer{60 - 40}\\
                      &= \answer{20}
    \end{align*}

\end{exe}


